\section{Tensor storage of Groundings}

\todo{
    A typical storage of a grounding is the enlistment of all non-zero coordinates.
    We understand this as a basis $\cpformat$ decomposition of the grounding tensor.
    We then propose more efficient storage tensor formats, which are respecting the structure of the grounded formula.
}

\subsection{Storage in Basis CP Decomposition}\label{sec:basisCPgrounding}

In many situations, grounding cores are sparse and representations as single tensor cores comes with a drastic overhead.
We often encounter sparse grounding tensors, where the number of non-zero coordinates (to be investigated by basis CP ranks in \charef{cha:sparseRepresentation}) satisfies
\begin{align*}
    \sparsityof{\groundingof{\folexformula}} << \inddim^{\cardof{\indvariableof{\folexformula}}} \, .
\end{align*}
In this case, since most coordinates vanish, the basis CP decomposition (see \secref{sec:basisCP}) enables a representation of the grounding with significantly lower storage demand, see \theref{the:sparseBasisCP}.
This is particularly useful for representing large relational databases, where each object has only a few relations with others, while the majority of possible relations remains unsatisfied.
We depict such CP decomposition of a formula grounding in \theref{fig:groundingCP}.

% Standard KB Encoding and Assumptions
Most logical syntaxes exploit $\ell_0$-sparsity, explicitly storing only known assertions.
The interpretation of unspecified assertions depends on the underlying assumptions.
Under the Closed World Assumption, all unspecified assertions are assumed to be false.

%
\begin{example}
    Consider the binary relation $Knows$ on the set of individuals $\{Alice,Bob,Charlie\}$ defined by the following pairs:
    \begin{center}

        \begin{tabular}{c|c}
            \textbf{Individual 0} & \textbf{Individual 1} \\
            \hline
            Alice                 & Bob                   \\
            Bob                   & Alice                 \\
            Charlie               & Bob                   \\
            Bob                   & Charlie
        \end{tabular}
    \end{center}

    \begin{align*}
        Knows[\indvariableof{0},\indvariableof{1}] =
        \coloredmatrixof{
            0 & 1 & 0 \\
            1 & 0 & 1 \\
            0 & 1 & 0
        }{\indvariableof{0},\indvariableof{1}} \, .
    \end{align*}
    with a $\cpformat$ representation with the leg cores
    \begin{align*}
        \legcoreofat{0}{\indvariableof{0},\decvariable}
        = \coloredmatrixof{
            1 & 0 & 0 & 0 \\
            0 & 1 & 1 & 0 \\
            0 & 0 & 0 & 1
        }{\indvariableof{0},\decvariable}
        \quad \text{and} \quad
        \legcoreofat{1}{\indvariableof{1},\decvariable}
        = \coloredmatrixof{
            0 & 1 & 0 & 0 \\
            1 & 0 & 0 & 1 \\
            0 & 0 & 1 & 0
        }{\indvariableof{1},\decvariable} \, .
    \end{align*}

    %%% For basis + CP
    Consider now an extension
    \begin{center}
        \begin{tabular}{c|c}
            \textbf{Individual 0} & \textbf{Individual 1} \\
            \hline
            Alice                 & Bob                   \\
            Bob                   & Alice                 \\
            Charlie               & Bob                   \\
            Bob                   & Charlie               \\
            John                  & Alice                 \\
            John                  & Bob                   \\
            John                  & John                  \\
            Alice                 & John                  \\
            Bob                   & John                  \\
            Charlie               & John                  \\
        \end{tabular}
    \end{center}
    Enumerating them in the order $[Alice,Bob,Charlie,John]$ this is the matrix
    \begin{align*}
        Knows[\indvariableof{0},\indvariableof{1}] =
        \coloredmatrixof{
            0 & 1 & 0 & 1\\
            1 & 0 & 1 & 1\\
            0 & 1 & 0 & 1\\
            1 & 1 & 1 & 1
        }{\indvariableof{0},\indvariableof{1}} \, .
    \end{align*}
    This has basis rank of $11$.
    Notice, that allowing for bas+ representation, we would reduce to $5$, by encoding that John knows everyone, and everyone knows John.

    \begin{align*}
        \legcoreofat{0}{\indvariableof{0},\decvariable}
        = \coloredmatrixof{
            1 & 0 & 0 & 0 & 0 & 1\\
            0 & 1 & 1 & 0 & 0 & 1\\
            0 & 0 & 0 & 1 & 0 & 1\\
            0 & 0 & 0 & 0 & 1 & 1
        }{\indvariableof{0},\decvariable}
        \quad \text{and} \quad
        \legcoreofat{1}{\indvariableof{1},\decvariable}
        = \coloredmatrixof{
            0 & 1 & 0 & 0 & 1 & 0\\
            1 & 0 & 0 & 1 & 1 & 0\\
            0 & 0 & 1 & 0 & 1 & 0\\
            0 & 0 & 0 & 0 & 1 & 1\\
        }{\indvariableof{1},\decvariable} \, .
    \end{align*}
\end{example}


\begin{figure}[t]
    \begin{center}
        \begin{tikzpicture}[scale=0.35, yscale=-1, thick] % , baseline = -3.5pt

    \draw (-1,1) rectangle (5,-1);
    \node[anchor=center] (text) at (2,0) {$\groundingof{\folexformula}$};

    \draw[] (0,1) -- (0,3) node[midway,left] {\colorlabelsize $\indvariableof{0}$};
    \draw[] (1.5,1) -- (1.5,3) node[midway,left] {\colorlabelsize $\indvariableof{1}$};
    \node[anchor=center] (text) at (2.75,2) {$\cdots$};
    \draw[] (4,1) -- (4,3) node[midway,right] {\colorlabelsize $\indvariableof{\variableorder\shortminus1}$};

    \node[anchor=center] (text) at (7,0) {${=}$};


    \begin{scope}
        [shift={(10,2)}]

        \coordinate (conposseldec) at (4.5,-5.5);

        \draw[fill] (conposseldec) circle (\dotsize);
        \draw (conposseldec) -- (4.5,-7.5) node[midway, right] {\colorlabelsize $\datvariable$};
        \draw[] (3.5,-7.5) rectangle (5.5, -9.5);
        \node[anchor=center] (text) at (4.5,-8.5) {\corelabelsize $\ones$};

        \draw[-<-] (0,1) -- (0,-1) node[midway,left] {\colorlabelsize $\indvariableof{0}$};
        \draw (-1,-1) rectangle (1, -3);
        \node[anchor=center] (text) at (0,-2) {\corelabelsize $\legcoreof{\folexformula,0}$};
        \draw[-<-] (0,-3) to[bend right=20] (conposseldec);


        \draw[-<-] (3,1) -- (3,-1) node[midway,left] {\colorlabelsize $\indvariableof{1}$};
        \draw (2,-1) rectangle (4, -3);
        \node[anchor=center] (text) at (3,-2) {\corelabelsize $\legcoreof{\folexformula,1}$};
        \draw[-<-] (3,-3) to[bend right=20]  (conposseldec);

        \node[anchor=center] (text) at (6,-2) {$\cdots$};

        \draw[-<-] (9,1) -- (9,-1) node[midway,left] {\colorlabelsize $\indvariableof{\variableorder\shortminus1}$};
        \draw (8,-1) rectangle (10, -3);
        \node[anchor=center] (text) at (9,-2) {\corelabelsize $\legcoreof{\folexformula,\variableorder\shortminus1}$};
        \draw[-<-] (9,-3) to[bend left=20]  (conposseldec);

    \end{scope}

\end{tikzpicture}
    \end{center}
    \caption{Basis CP Decomposition of the grounding of $\folexformula$, following the scheme of \theref{the:sparseBasisCP}.
    Instead of direct storage of the grounding tensor $\groundingof{\folexformula}$, the non-zero coordinates are enumerated by a variable $\datvariable$ and the corresponding coordinates stored in leg-matrices $\legcoreof{\folexformula,\indenumerator}$.}
    \label{fig:groundingCP}
\end{figure}



\subsection{Storage in Tensor Network Formats}

We now investigate, whether the basis $\cpformat$ storage of a formula is optimal.



\begin{itemize}
    \item Advantage of contracted storage: Redundancies by vanishing slices removed.
    We could try to remove redundancies by message passing algorithms, processing constraints on the support of the tensor.
    \item Advantage of decomposed store: Repetition effects of partial indices are respected.
\end{itemize}


\begin{example}[Conjunction chain of random predicates]

    Let us consider a conjunction chain of predicates, i.e. the $\catenumerator$ predicate shares with the $\catenumerator+1$ and $\catenumerator-1$ predicate one term variable, but with no other predicate.
    The formula would be a $\ttformat$-like tensor network.
    %
    Now let us consider an random world, where each complete substitution of each predicate is satisfied independently with probability $p$.
    The expectation of the groundings is the tensor $p^{\catorder} \cdot \onesat{\indvariableof{[\catorder+1]}}$ and the expectation of the contraction is
    \begin{align*}
        p^{\catorder} \cdot \inddim^(\catorder+1) \, .
    \end{align*}
    The expected basis $\cpformat$ rank of the grounding grows exponentially in $\catorder$, while in the $\ttformat$-like decomposition the expectation of the storage demand is
    \begin{align*}
        p \cdot \catorder \cdot \inddim^2 \,
    \end{align*}
    and thus linear in $\catorder$.

    % Arbitrary random tensor network
    More general, when any random network of predicates is considered, the expected contraction is
    \begin{align*}
        p^{\cardof{\edges}} \cdot \inddim^{\indorder}
    \end{align*}
    and the storage demand
    \begin{align*}
        \sum_{\edgein} p \cdot \inddim^{\cardof{\edge}} \, .
    \end{align*}



\end{example}


\subsection{Example: Relational Databases}

A database is understood as a specific \firstOrderLogic{} world, and are operations on such a single world.
Queries are described by a formula $\impformula$, which are asked against a specific world $\worldindices$ to retrieve the grounding $\groundingof{\impformula}$.
The variables of such formulas are called projection variables.
The answer $\groundingof{\impformula}$ of a query is most conveniently represented as a list of solution mappings from the projection variables to objects in the world, such that the query formula is satisfied.
Answering a query by solution mappings corresponds with finding the basis CP Decomposition (see \secref{sec:basisCP}) of $\groundingof{\impformula}$.
We can understand these solution mappings as stored in the leg-matrices $\legcoreof{\folexformula,\indenumerator}$ (see \figref{fig:groundingCP}).

Let us give with the outer join an example of a popular operation to define queries, which efficient execution and storage can be improved based on considerations in the tensor network formalism.

\begin{definition}[Outer join]
    Let there be a world $\worldindices$ and formulas $\extformulaof{\selindex}$ depending on variables $\indvariableof{\nodesof{\selindex}}$, which have grounding tensors by
    \begin{align*}
        \groundingofat{\extformulaof{\selindex}}{\indvariableof{\node}} \wcols  \bigtimes_{\node\in\nodesof{\selindex}}[\inddimof{\node}] \rightarrow \ozset \, .
    \end{align*}
    Then their (outer) $\joinsymbol$ is defined as the grounding of their conjunctions, as
    \begin{align*}
        \groundingofat{\joinsymbol\left(\extformulaof{0},\ldots,\extformulaof{\seldim-1}\right)}{\bigcup_{\selindexin}\indvariableof{\nodesof{\selindex}}}
        = \contractionof{\groundingofat{\extformulaof{\selindex}}{\indvariableof{\nodesof{\selindex}}}\,:\,\selindexin}{\bigcup_{l\in[p]}\indvariableof{\nodesof{\selindex}}} \, .
    \end{align*}
\end{definition}

%Visualization and efficiency
We can understand the $\joinsymbol$ of groundings by a factor graph, where each grounding tensor decorates the hyperedge to the node set $\nodesof{\selindex}$.
The projection variable assignment to each formula combined in a $\joinsymbol$ operation provide a basic tensor network format to store the output of the operation.
There are thus situations, in which the solution map storage corresponding with a CP Decomposition comes with unnecessary overheads compared with other formats.

% Coordinatewise transform
We can also understand the $\joinsymbol$ operation as a coordinatewise transform (see \defref{def:coordinatewiseTransform}) with the product as transform function.
To make this connection solid, one would need to extend each joined formula trivially to the variables appearing in other formulas.

% Evaluation similar constraint propagation
The efficiency of evaluating the contraction to a $\joinsymbol$ operation might be improved by understanding it as an Constraint Satisfaction Problem (see \charef{cha:logicalReasoning}).
When applying efficient Message Passing algorithms such as Knowledge Propagation (see \algoref{alg:knowledgePropagation}), the groundings can be sparsified by local constraint propagation operations before turning to more global and more demanding contraction operations.
Here the groundings $\groundingof{\extformulaof{\selindex}}$ would be used to initialize Knowledge Cores $\kcoreof{\edge}$ and sequentially sparsified during the algorithm.

%\begin{example} % WOULD NEED OVERWORK: DRAW!
%	For example take a query with many basic graph patterns with pairwise different projection variables.
%	The global CP Decomposition would come here with an exponential storage overhead compared with storage as a tensor product of CP Decompositions to each Basic graph pattern.
%\end{example}

%% CONFUSING?
%\begin{remark}[Distinguishing from probabilistic queries]
%	Let us distinguish the discussion here from those of queries in probabilistic reasoning, which have two main differences.
%	First, we ask queries against all possible pairs of variables, instead of asking the probability of satisfaction of a specific formula.
%	Second, since we made the epistemologic assumption of knowing possibilities and not probabilities in logics, a query is answered by a truth value.
%	We then only output in the shape of solution mappings the variable assignments where the query formula is true.
% 	Thus, the queries here can be thought of as a batch of probabilistic queries with Boolean answers.
%	% Alternative -> Later?
%	Probabilistic queries can furthermore be understood in terms of the data extraction process described in this section.
%	We can ask the query in probabilistic form (decomposed into atomic formulas) on the resulting empirical distribution.
%	This results in the ratio of the worlds satisfying the query among those worlds satisfying the extraction query $\impformula$.
%\end{remark}

